% !PRINTED 469
\weberpage{469}
%<<CONT

also ergibt (10):
\weq{12}{B\left(m\,\frac{A+C}{2} - (m+1)\,B + A\right) \equiv 0 \quad (\mathrm{mod}\ 8).}
Wenn nun $m \equiv 3\ (\mathrm{mod}\ 4)$, so ist $A + C \equiv 0\ (\mathrm{mod}\ 4)$ und der in (12) in der Klammer stehende Ausdruck ungerade, daher ist (12) \gesperrt{nur} unter der Voraussetzung befriedigt, daß
\[
B \equiv 0 \quad (\mathrm{mod}\ 8).
\]
Wir wollen, um Wiederholungen zu vermeiden, bei den im folgenden auszusprechenden Theoremen ein- für allemal voraussetzen, daß $\omega$ die Wurzel einer solchen Gleichung (3) der Diskriminante $-4m$ sei, in der $A$ ungerade, $B$ durch 8 teilbar sei. Damit ist also das Theorem bewiesen:

4. \gesperrt{Ist $m \equiv 3\ (\mathrm{mod}\ 4)$, so ist $f(\omega)^6$ Klasseninvariante.}

Ist ferner $m \equiv 5\ (\mathrm{mod}\ 8)$, so reduziert sich die Kongruenz (12) auf:
\[
B\left(\frac{5C - A}{2} + 2B\right) \equiv 0 \quad (\mathrm{mod}\ 8),
\]
und diese Bedingung ist erfüllt, wenn $B$ durch 4 teilbar ist, dagegen nicht erfüllt, wenn $B$ nur durch 2, nicht durch 4 teilbar ist. Daraus folgt:

5. \gesperrt{Ist $m \equiv 5\ (\mathrm{mod}\ 8)$, so ist $f(\omega)^{12}$ Klasseninvariante.}

Ist $m \equiv 1\ (\mathrm{mod}\ 8)$, so läßt sich aus der Kongruenz (12) nichts schließen. Dies stimmt mit dem Umstande überein, daß in diesem Falle $\Phi_m(u, u)$ nur von $u^8$ abhängig ist. Um auch hier zu ähnlichen Resultaten zu gelangen, wenden wir die Transformation zweiter Ordnung an. Wir nehmen in (7):
\weq{13}{n \equiv 1 \quad (\mathrm{mod}\ 8)}
und setzen:
\[
u\,v = \pm \sqrt{2}^{3},
\]
d. h.:
\[
f\left(\frac{c + \partial\omega}{a}\right)^{3} f(\omega)^{3} = \pm \sqrt{2}^{3},
\]
oder [§ 34, (18)]:
\weq{14}{f\left(\frac{c + \partial\omega}{a}\right)^{3} = \pm f\left(\frac{\omega - 1}{\omega + 1}\right)^{3}.}
Die Gleichung (7) geht dadurch über in eine Gleichung:
\weq{15}{\Phi_n\left(u, \frac{\pm\sqrt{2}^{3}}{u}\right) = 0,}
%CONT>>
